\begin{theorem}[Punctured-domain many-one reduction]
\label{thm:punctured-domain-richardson-reduction}
\label{thm:richardson-punctured-domain-convention}
Fix finite pure EML terms \((\mathsf P,\mathsf S,\mathsf A)\) satisfying
\[
  \mathsf P=\pi,
  \qquad \mathsf S(x)=\sin x\quad(x\in\RR),
  \qquad \mathsf A(x)=|x|\quad(x\ne0).
\]
For each Richardson expression \(E\), define its punctured Richardson domain
\(D_E\subseteq\RR\) recursively as the set of real inputs for which every
ordinary elementary subterm is defined and, whenever a subterm \(\mu(F)\) is
evaluated, the value of \(F\) is nonzero.  Equivalently, the transported EML
instance is the source-indexed partial function obtained by restricting the
translated term to this same recursively defined set \(D_E\), not by enlarging it
to any principal-branch domain on which the translated syntax happens to have a
value.  Then
\(\tau_{\mathsf P,\mathsf S,\mathsf A}(E)\) and
\(\tau_{\mathsf P,\mathsf S,\mathsf A}(0)\) are defined on \(D_E\), the first term
equals \(E\) there, and the second term equals \(0\) there.  In particular,
\[
  E=0\hbox{ on }D_E
  \quad\Longleftrightarrow\quad
  \tau_{\mathsf P,\mathsf S,\mathsf A}(E)=
  \tau_{\mathsf P,\mathsf S,\mathsf A}(0)
  \hbox{ on }D_E .
\]
Zeros of absolute-value arguments therefore delete the same input points from the
source-indexed transported equality problem on both sides and do not invalidate
the many-one reduction.  The endpoint row \(U_{126}\) is not used in this
punctured obstruction: replacing \(\mathsf A\) by the internally proved
absolute-value row \(\Abs\), or by any separately verified punctured
absolute-value row, is the whole root interface required here.
\end{theorem}

\begin{proof}
Proceed by structural induction on \(E\), proving simultaneously that the source
subterm and its translated EML term agree on the recursively transported domain
and that the transported domain is the one assigned by the same recursion.
Constants, the variable, arithmetic, composition, and \(\exp\) use the internally
proved core identities and their ordinary real domains.  The \(\pi\) and \(\sin\)
clauses use the first two displayed hypotheses.  For a subterm \(\mu(F)\), the
recursive definition of \(D_{\mu(F)}\) requires \(x\in D_F\) and \(F(x)\ne0\).  By
the induction
hypothesis,
\(\tau_{\mathsf P,\mathsf S,\mathsf A}(F)(x)=F(x)\), so the translated argument to
\(\mathsf A\) is nonzero exactly on the same points.  The third displayed
hypothesis gives
\[
  \mathsf A(\tau_{\mathsf P,\mathsf S,\mathsf A}(F)(x))
  =|F(x)|=
  \mu(F)(x)
\]
for every \(x\in D_{\mu(F)}\).  If \(F(x)=0\), the recursive convention removes
\(x\) from the punctured domain before either the Richardson \(\mu\)-clause or the
EML \(\mathsf A\)-clause is used; no value at that point is needed, and no
endpoint-root semantics can repair or change the deleted point.  Thus each
constructor preserves both the relevant domain and the value on that domain.
Taking \(E=0\) gives the displayed equivalence, which is an effective many-one
reduction because the translation is computable.
\end{proof}
\begin{theorem}[Finite-budget Richardson translation]
\label{thm:richardson-translation-budget}
\label{lem:richardson-translation-bookkeeping}
Fix finite pure EML replacement terms \((\mathsf P,\mathsf S,\mathsf A)\) and a
finite EML witness table for the arithmetic, exponential, rational-constant,
\(\log2\), \(\pi\), sine, and punctured absolute-value rows used in
\Cref{lem:richardson-translation}.  Let \(m_R\) be the number of pure EML nodes
in the replacement macro for a Richardson constructor \(R\), counted after
expanding all internal arithmetic macros and counting each argument placeholder
once.  Let \(d_R\) be the number of argument-placeholder occurrences in that
macro, and set \(B=\max_R m_R\) and \(d=\max_R d_R\).
Then \(B\) and \(d\) are finite and effectively computable from the printed
witness table.  If a Richardson expression \(E\) has \(n\) source-constructor
nodes and no source subexpression is shared, then the fully expanded pure EML
syntax of \(\tau_{\mathsf P,\mathsf S,\mathsf A}(E)\) has at most
\[
  B(1+d+d^2+\cdots+d^n)\le B(n+1)\max(1,d)^n
\]
pure EML nodes.  Consequently the translation is a uniform finite-substitution
procedure: every output term is obtained by at most this many explicit node
insertions, the bound depends only on the fixed witness table and \(n\), and the
punctured-domain equality reduction of
\Cref{thm:punctured-domain-richardson-reduction} is a genuine finite many-one
reduction rather than an appeal to an unexpanded compiler.
\end{theorem}

\begin{proof}
There are only finitely many source constructors and, by hypothesis, each row in
the witness table is a printed finite pure EML macro after expanding the internal
arithmetic abbreviations.  Hence the two maxima defining \(B\) and \(d\) are
finite; they are computable by parsing the finite strings and counting pure EML
nodes and placeholder occurrences.

We prove the node bound by induction on the source tree.  Let \(N(F)\) be the
number of pure EML nodes in the fully expanded translation of a source
subexpression \(F\), and let \(h(F)\) be the number of source-constructor nodes in
\(F\).  If \(F\) is obtained by applying a constructor \(R\) to immediate
subexpressions \(F_1,\ldots,F_r\), then full expansion substitutes a separate
copy of \(\tau(F_i)\) into each placeholder occurrence assigned to the \(i\)-th
argument.  If \(q_{R,i}\) denotes the number of such occurrences, then
\(\sum_i q_{R,i}\le d\), and therefore
\[
  N(F)\le B+\sum_i q_{R,i}N(F_i)
       \le B+d\max_i N(F_i)
\]
with the evident convention for nullary constructors.  Since each proper
subexpression has height at most \(h(F)-1\), iteration along any root-to-leaf path
gives
\[
  N(F)\le B(1+d+d^2+\cdots+d^{h(F)}) .
\]
Because the height of a tree with \(n\) constructor nodes is at most \(n\), this
implies the displayed estimate for \(E\).  The recursive construction also gives
an explicit finite substitution algorithm: visit the source tree from the
leaves upward, replace the constructor at each node by its fixed macro, and copy
the already constructed child strings into the finitely many placeholders.  The
algorithm performs no semantic search and uses no external catalogue beyond the
fixed finite witness table.  Combining this finite construction with the
domain-preservation theorem \Cref{thm:punctured-domain-richardson-reduction}
therefore yields the stated finite many-one reduction.
\end{proof}

\begin{theorem}[Witness-row rigidity and defect classification]
\label{thm:richardson-witness-row-rigidity}
\label{thm:richardson-row-necessity-and-sufficiency}
Let \((\mathsf P,\mathsf S,\mathsf A)\) be finite pure EML terms.  The structural
translation \(\tau_{\mathsf P,\mathsf S,\mathsf A}\), with \(\mathsf P\),
\(\mathsf S\), and \(\mathsf A\) substituted respectively for Richardson's
\(\pi\), \(\sin\), and punctured absolute-value primitive, gives a
value-preserving many-one reduction from Richardson's punctured
identity-to-zero problem to a source-indexed pure EML equality problem if and
only if the three witness rows
\[
  \mathsf P=\pi,
  \qquad \mathsf S(x)=\sin x\quad(x\in\RR),
  \qquad \mathsf A(x)=|x|\quad(x\ne0)
\]
hold in the principal-branch EML semantics.  Consequently the exact printed
pair \((U_{47},U_{111})\) is not merely missing a parser certificate: the two
real-domain identities for those exact strings, either as the explicit
\((H_{\mathcal W})\) hypothesis or as separately proved facts, are logically
necessary and sufficient for the printed-pair specialization of the Richardson
obstruction, once the internally proved punctured row \(\mathsf A=\Abs\) is used.

Equivalently, the only possible failures of a finite source-locked Richardson
package are the three row defects
\[
  \mathsf P\ne\pi,
  \qquad \mathsf S\not\equiv\sin\hbox{ on }\RR,
  \qquad \mathsf A\not\equiv |-|\hbox{ on }\RR\setminus\{0\}.
\]
Here a row defect means that the displayed term is either undefined at some
required real input or defined there with the wrong value.  If any one of these
defects occurs, the structural translation already fails to be value-preserving
on the corresponding primitive test expression
  \(\pi\), \(\sin(x)\), or \(\mu(x)\).  Equivalently, the bad examples for this
  interface form the three classifiable row-defect skeletons: a wrong constant
  row, a wrong unary sine row, or a wrong punctured absolute-value row.  Hence no
  equality algorithm, semantic
normalizer, or protocol compiler obstruction in this paper may be attributed to
that defective finite triple.  If none occurs, \Cref{thm:separation} gives the
full punctured Richardson obstruction for that triple.
\end{theorem}

\begin{proof}
If the three displayed rows hold, \Cref{thm:punctured-domain-richardson-reduction}
gives value preservation on every punctured Richardson domain and hence the
many-one reduction to the transported EML equality problem.

Conversely suppose the indicated structural translation is value-preserving for
Richardson expressions on their punctured domains.  Apply it first to the
constant expression \(\pi\).  Its domain is all of \(\RR\), and the translated
term is \(\mathsf P\); value preservation gives \(\mathsf P=\pi\).  Apply it next
to the expression \(\sin(x)\).  Its domain is again all of \(\RR\), and the
translated term is \(\mathsf S(x)\); value preservation gives
\(\mathsf S(x)=\sin x\) for every real \(x\).  Finally fix an arbitrary nonzero
  real number \(a\) and apply the translation to the one-variable Richardson
  expression \(\mu(x)\) at the input value \(x=a\).  The punctured domain of
  \(\mu(x)\) is \(\RR\setminus\{0\}\), so \(a\) is an allowed input, the translated value is
\(\mathsf A(a)\), and Richardson's primitive has value \(|a|\).  Thus
\(\mathsf A(a)=|a|\) for every \(a\ne0\).  These three test expressions are
primitive clauses of the same structural recursion used in every larger
Richardson expression, so no weaker row data can supply the stated
value-preserving reduction.

Taking \(\mathsf A=\Abs\), the third row is supplied by
\Cref{thm:eml-punctured-richardson-core}.  Therefore the printed-pair
specialization with \((U_{47},U_{111})\) exists exactly when the first two tests
above hold for those printed terms; finite syntax alone does not enter the
argument except to make the translated EML terms finite.

The defect classification is the contrapositive of the three primitive tests
just used, together with the sufficiency direction.  If the \(\pi\)-row is
defective, value preservation fails on the constant Richardson expression
\(\pi\).  If the sine row is defective, choose a real input where
\(\mathsf S\) is undefined or has value different from \(\sin\); value
preservation fails for the primitive expression \(\sin(x)\) at that input.  If
the absolute-value row is defective, choose a nonzero real input where
\(\mathsf A\) is undefined or has value different from \(|-|\); value
preservation fails for the primitive punctured expression \(\mu(x)\) there.  In
each defective case the many-one reduction is not the Richardson reduction
proved in this manuscript, so the downstream zero-test, normalizer, and compiler
obstructions cannot be invoked for that finite triple.  In the no-defect case
the three displayed identities hold, and \Cref{thm:separation} supplies exactly
those downstream obstructions.
\end{proof}

\begin{theorem}[Primary verified replacement-triple Richardson obstruction]
\label{thm:proc-ams-relative-richardson-theorem}
\label{thm:relative-richardson-main-spine-closure}
\label{thm:interface-only-richardson-spine}
\label{thm:printed-root-richardson-conditional-clause}
\label{thm:no-unconditional-printed-richardson-edge}
\label{thm:separation}
\label{thm:no-principal-branch-eml-nonnegative-root}
Fix the principal-branch pure EML semantics of \Cref{def:eml}.  The primary,
unconditional Richardson theorem used by this paper is the following
replacement-triple scheme: for any finite pure EML triple
\((\mathsf P,\mathsf S,\mathsf A)\) satisfying
\[
  \mathsf P=\pi,
  \qquad \mathsf S(x)=\sin x\quad(x\in\RR),
  \qquad \mathsf A(x)=|x|\quad(x\ne0).
\]
Richardson's punctured identity-to-zero problem many-one reduces to the
source-indexed partial equality problem for the pure EML fragment translated
using this verified triple; the source domain \(D_E\) is part of the instance and
equality is equality of the transported partial real functions on that domain.
Consequently that transported fragment admits no zero-identity algorithm, no
computable semantic normalizer into finite objects with decidable equality that
reflects and preserves equality on the transported source domains, and no
equality-reflecting semantic compiler into finite decidable protocol objects.

The theorem has an unconditional scheme instance for every triple once the three
displayed real-domain identities for that triple are available as proved
hypotheses.  In particular the internally proved punctured row \(\Abs(u)=|u|\) from
\Cref{thm:eml-punctured-richardson-core} may serve as \(\mathsf A\).  The exact
printed pair \((U_{47},U_{111})\) gives the printed-string specialization of this
theorem only after the two real-domain identities
\(U_{47}=\pi\) and \(U_{111}(x)=\sin x\) for \(x\in\RR\) have been supplied in
the principal-branch semantics.  Equivalently, a citation to this theorem with
the printed pair is shorthand for the optional hypothesis \((H_{\mathcal W})\)
or for independently verified replacements of exactly those two identities.  The
finite syntax catalogue by itself is not one of the verified triples covered by
the theorem.  Thus every submission-facing use of the printed pair is a
conditional branch, whereas every submission-facing unconditional Richardson
conclusion factors through the abstract verified replacement triple above.
Without \((H_{\mathcal W})\), or without separate analytic proofs of those exact
two printed rows, the strings \(U_{47}\) and \(U_{111}\) do not occur in the
unconditional conclusion of this theorem.

Moreover, no finite pure EML term defined on a right-neighbourhood of \(0\) can
represent \(t\mapsto\sqrt t\) on \(\RR_{\ge0}\).  This endpoint obstruction does
not affect the preceding Richardson reduction, because that reduction uses only
the punctured absolute-value identity \(\mathsf A(u)=|u|\) for \(u\ne0\).  The
multiple labels on this theorem are dependency-cut aliases for the same verified
triple interface and its printed-string boundary.
\end{theorem}

\begin{proof}
By \Cref{lem:richardson-translation,thm:richardson-translation-budget}, the
translation \(\tau_{\mathsf P,\mathsf S,\mathsf A}\) is effective.  By
\Cref{thm:punctured-domain-richardson-reduction}, for every Richardson expression
\(E\),
\[
  E=0\hbox{ on }D_E
  \quad\Longleftrightarrow\quad
  \tau_{\mathsf P,\mathsf S,\mathsf A}(E)=
  \tau_{\mathsf P,\mathsf S,\mathsf A}(0)
  \hbox{ on }D_E .
\]
Since Richardson's punctured identity-to-zero problem for
\(\mathcal R_{\mathrm{core}}\) is undecidable by
\Cref{thm:richardson-punctured-source-problem}, this is a many-one reduction.  A
zero-identity algorithm on the source-indexed translated fragment would decide
the Richardson problem.  A computable semantic normalizer into finite objects
with decidable equality would decide the same problem only if it reflects and
preserves equality of the source-indexed partial real functions: normalize the
restrictions of \(\tau_{\mathsf P,\mathsf S,\mathsf A}(E)\) and
\(\tau_{\mathsf P,\mathsf S,\mathsf A}(0)\) to the same domain \(D_E\) and compare
the results.  This is the transported-normalizer obstruction of
\Cref{lem:transported-partial-normalizer-obstruction}.  Likewise, an
equality-reflecting semantic compiler into any finite decidable protocol object
would decide the problem by compiling the two translated terms and checking
target equality.

The specialization statement is the rigidity assertion of
\Cref{thm:richardson-witness-row-rigidity} applied to the printed pair.  The
preceding argument used only the three displayed identities for the chosen
triple \((\mathsf P,\mathsf S,\mathsf A)\), so any triple for which those
identities are hypotheses of the theorem is an instance of the interface.  The term
\(\Abs=G(M(-,-))\) satisfies the third identity by
\Cref{thm:eml-punctured-richardson-core}.  For the printed names
\(U_{47}\) and \(U_{111}\), however,
\Cref{prop:printed-pi-sine-instantiation-criterion,cor:no-unconditional-printed-richardson-triple}
prove exactly that the catalogue supplies finite syntax but not the two semantic
base identities for \(\pi\) and sine.  Thus those printed strings are not part of
the unconditional theorem scheme; they enter only under the corresponding
verified-row hypothesis, namely the \(U_{47},U_{111}\) part of
\((H_{\mathcal W})\), or under separate proofs of the same two identities.

It remains only to record why this does not smuggle in an endpoint square-root
term.  Let \(T(t)\) be a finite pure EML term defined for all
\(t\in[0,\epsilon)\).  After possibly shrinking \(\epsilon\), the real-domain
value of every subterm of \(T\) is real-analytic near \(0\): constants and the
input are analytic, and if a subterm is \(\eml(A(t),B(t))\), then definedness at
\(0\) gives \(B(0)\ne0\), so \(B\) stays in a single principal-logarithm branch
neighbourhood and \(e^{A(t)}-\operatorname{Log}B(t)\) is analytic.  Hence a term
equal to \(\sqrt t\) on \([0,\epsilon)\) would be analytic at \(0\).  Its Taylor
series would have the form \(a_0+a_1t+a_2t^2+\cdots\); equality with \(\sqrt t\)
for \(t>0\) forces \(a_0=0\), and division by \(t\) gives a bounded analytic
left side equal to \(t^{-1/2}\), a contradiction as \(t\downarrow0\).  Thus the
endpoint root is impossible, while the punctured term \(G(M(u,u))\) remains
valid for every nonzero real argument and is exactly the condition Richardson
uses.
\end{proof}

\begin{proposition}[Three-witness stability of the Richardson interface]
\label{prop:three-witness-stability}
\label{prop:three-witness-richardson-stability}
Let \((\mathsf P,\mathsf S,\mathsf A)\) be finite pure EML terms satisfying the
three real-domain identities
\[
  \mathsf P=\pi,
  \qquad \mathsf S(x)=\sin x\quad(x\in\RR),
  \qquad \mathsf A(x)=|x|\quad(x\ne0).
\]
Replacing the displayed witnesses \((W_{\pi},W_{\sin},\Abs)\) by
\((\mathsf P,\mathsf S,\mathsf A)\) gives an effective translated Richardson
fragment with exactly the obstruction asserted in \Cref{thm:separation}: no
zero-identity algorithm, no computable semantic normalizer into finite decidable
objects that reflects and preserves equality on the transported source domains,
and no equality-reflecting semantic compiler into finite decidable protocol
objects.  Hence the connection from \Cref{thm:separation} to
\Cref{cor:strict-L0-L1,thm:obstruction-spine,thm:main-hierarchy} depends only on
these three identities, not on the printed syntax of the chosen witnesses.  The
printed pair \((U_{47},U_{111})\) instantiates the first two rows only after those
two identities have been proved for the exact strings; the third row in the
submission theorem is the punctured term \(\Abs\), not \(U_{126}\).
\end{proposition}

\begin{proof}
Use the translation \(\tau_{\mathsf P,\mathsf S,\mathsf A}\) of
\Cref{lem:richardson-translation}, using
the unchanged clauses for rational constants, \(\log2\), the variable,
\(+,-,\times\), composition, and \(\exp\), and using \(\mathsf P\),
\(\mathsf S\), and \(\mathsf A\) in the clauses for \(\pi\), \(\sin\), and the
Richardson unary symbol \(\mu\).  The recursion is effective because the source
expression tree and the three replacement terms are finite.  An induction on
Richardson expressions gives value preservation on the real domains relevant to
Richardson's theorem: all unchanged clauses are those already verified in
\Cref{lem:eml-core-witness-table}; the three changed clauses are exactly the
three displayed hypotheses; and the only partial clause is the \(\mu\)-clause,
where Richardson requires agreement with \(|x|\) only for nonzero arguments.

Thus, for every expression \(E\) in \(\mathcal R_{\mathrm{core}}\),
\Cref{thm:punctured-domain-richardson-reduction} gives the equivalence of
identity to zero on the punctured Richardson domain and identity to the
translated zero term on the same domain.  Since the punctured identity-to-zero
problem for \(\mathcal R_{\mathrm{core}}\) is undecidable by
\Cref{thm:richardson-punctured-source-problem}, this is a many-one reduction to
the source-indexed replacement EML fragment.  A zero-identity algorithm for that
fragment would decide Richardson identity-to-zero; a computable semantic
normalizer into finite objects with decidable equality would decide it only under
the stated equality-reflection and preservation condition on the
source-restricted partial real functions, by normalizing the two source-restricted
translated terms and comparing the normal forms; and an equality-reflecting
semantic compiler into finite decidable protocol objects would decide it by
compiling the two source-restricted translated terms and comparing the target
objects.  These are precisely the
obstruction consequences used downstream in
\Cref{cor:strict-L0-L1,thm:obstruction-spine,thm:main-hierarchy}.  Finally, the
argument has used only the three displayed semantic identities.  Therefore a
specific printed pair, including \((U_{47},U_{111})\), supplies the first two rows
exactly when those identities have been proved for that pair; its name or
provenance supplies no additional semantic content.  The third row in the
submission theorem is supplied by the punctured term \(\Abs\), unless a different
finite punctured absolute-value witness is separately proved.
\end{proof}
The submission Richardson scope is therefore the following.  The Proc. AMS core
Richardson theorem stream is unconditional as a scheme over verified finite
punctured triples \((\mathsf P,\mathsf S,\mathsf A)\): once the three displayed
identities for \(\pi\), sine, and \(|x|\) off zero have been proved for a finite
triple, the obstruction follows.  The paper proves the internal absolute-value
option \(\mathsf A=\Abs\); the printed pair \((W_{\pi},W_{\sin})\) instantiates the
other two rows only under \((H_{\mathcal W})\).  It contains no unconditional
endpoint-root instance for \(U_{126}\).  This is an immediate extraction from
\Cref{thm:separation,prop:three-witness-stability}: the endpoint obstruction in
\Cref{thm:separation} shows why no pure EML term defined at \(0\) can be
substituted for the positive root term, so the Richardson input to
\Cref{thm:obstruction-spine,thm:main-hierarchy} is the punctured triple
interface, while the endpoint-root string \(U_{126}\) remains outside that input.

\begin{theorem}[Maximal verified-triple closure of the Richardson clause]
\label{thm:unconditional-submission-richardson-closure}
\label{thm:verified-triple-only-main-richardson-clause}
In the submission manuscript with no appeal to \((H_{\mathcal W})\), the
Richardson clause used by
\Cref{thm:obstruction-spine,thm:proc-ams-core-obstruction-spine,thm:main-hierarchy}
is exactly the following maximal verified-triple theorem and nothing stronger;
this is the form to which every submission-facing Richardson citation is to be
read:
for every finite pure EML triple \((\mathsf P,\mathsf S,\mathsf A)\) satisfying
\[
  \mathsf P=\pi,
  \qquad \mathsf S(x)=\sin x\quad(x\in\RR),
  \qquad \mathsf A(x)=|x|\quad(x\ne0),
\]
the source-indexed EML fragment obtained from Richardson's punctured language by
\(\tau_{\mathsf P,\mathsf S,\mathsf A}\) has no zero-identity algorithm, no
computable semantic normalizer into finite decidable objects that reflects and
preserves equality on the transported source domains, and no equality-reflecting
semantic compiler into finite decidable protocol objects.
Conversely, if a source-locked finite pure EML triple gives the same
value-preserving Richardson translation and hence the same three obstruction
consequences through the structural recursion \(\tau\), then the three displayed
rows must hold.  Thus the verified-triple condition is not merely a sufficient
hypothesis but the necessary and sufficient interface for this paper's
Richardson theorem chain.
The manuscript proves the third row unconditionally for
\(\mathsf A=\Abs=G(M(-,-))\), but it does not prove the first two rows for the
particular printed names \((U_{47},U_{111})\).  Hence the exact printed pair
\((U_{47},U_{111})\) is not an unconditional instance of the displayed theorem;
it becomes an instance precisely after those two printed real-domain identities
are assumed as \((H_{\mathcal W})\) or supplied as proved hypotheses.  More generally,
\Cref{thm:richardson-witness-row-rigidity} classifies every source-locked finite
triple by its row defects: a defective \(\pi\), sine, or punctured absolute-value
row fails already on a primitive Richardson test expression, while a defect-free
triple is exactly a verified triple covered by \Cref{thm:separation}.  The cited
\(\mathcal B_{36}\) representative existence statement and the finite syntax
catalogue do not alter this closure; by the interface rigidity of
\Cref{thm:submission-interface-rigidity}, deleting \((H_{\mathcal W})\) removes
exactly the named printed-pair branch and no replacement-triple consequence of
the main hierarchy.
\end{theorem}

\begin{proof}
The displayed verified-triple statement is exactly \Cref{thm:separation}.  Its
proof uses Richardson's punctured undecidability theorem through
\Cref{thm:richardson-punctured-source-problem}, the effective translation and
domain preservation of
\Cref{lem:richardson-translation,thm:punctured-domain-richardson-reduction}, and
the three displayed semantic identities.  The normalizer and compiler exclusions
are the finite-decidable-equality consequences proved in the same theorem and in
\Cref{prop:L1-compiles-L0}.  Thus every downstream citation from
\Cref{thm:obstruction-spine,thm:proc-ams-core-obstruction-spine,thm:main-hierarchy}
to the Richardson obstruction factors through this verified-triple interface.

The internal proof of the third row for \(\Abs\) is
\Cref{thm:eml-punctured-richardson-core}: for nonzero real \(u\),
\(M(u,u)=u^2>0\) and the positive-domain square-root witness gives
\(G(M(u,u))=\sqrt{u^2}=|u|\).  The first two rows for the printed names are
different.  By \Cref{lem:eml-pi-sine-witnesses} the catalogue proves only that
\(U_{47}\) and \(U_{111}\) are finite pure EML strings; by
\Cref{prop:printed-pi-sine-instantiation-criterion,cor:no-unconditional-printed-richardson-triple}
the identities \(U_{47}=\pi\) and \(U_{111}(x)=\sin x\) on \(\RR\) are exactly the
extra hypotheses needed to substitute those names for \((\mathsf P,\mathsf S)\).
Therefore omitting \((H_{\mathcal W})\) removes only the printed-name
specialization, not the verified-triple theorem itself.  The stronger row-defect
classification cited in the statement is \Cref{thm:richardson-witness-row-rigidity}:
if a source-locked finite triple is not a verified triple, one of its three
primitive rows fails and the structural Richardson reduction is unavailable
already at the corresponding primitive test expression.

Finally,
\Cref{thm:B36-import-exact-scope,thm:B36-self-contained-replacement-core,thm:submission-interface-rigidity}
show that the \(\mathcal B_{36}\) import is representative existence for the
external calculator basis and is not a term-by-term proof of the Richardson
witness rows.  The interface-rigidity theorem also identifies the exact effect
of removing \((H_{\mathcal W})\): it deletes the named printed-pair branch while
leaving the replacement-triple Richardson obstruction used by
\Cref{thm:main-hierarchy}.  Consequently neither that citation nor the syntax
catalogue can create an unconditional printed \((U_{47},U_{111})\) Richardson
instance, and the submission-facing Richardson clause is precisely the
verified-triple closure stated above.
\end{proof}

\begin{corollary}[Printed-string-free Richardson spine]
\label{cor:printed-string-free-richardson-spine}
If every occurrence of the optional hypothesis \((H_{\mathcal W})\) and every
unproved printed identification of \((U_{47},U_{111})\) with \((\pi,\sin)\) is
removed from the submission, the Richardson part of
\Cref{thm:obstruction-spine,thm:proc-ams-core-obstruction-spine,thm:main-hierarchy}
remains the following unconditional theorem scheme: for each finite pure EML
triple \((\mathsf P,\mathsf S,\mathsf A)\) satisfying
\[
  \mathsf P=\pi,
  \qquad \mathsf S(x)=\sin x\quad(x\in\RR),
  \qquad \mathsf A(x)=|x|\quad(x\ne0),
\]
the transported punctured Richardson fragment has no zero-identity algorithm, no
computable semantic normalizer into finite decidable objects reflecting and
preserving equality on source domains, and no equality-reflecting compiler into
finite decidable protocol objects.  The optional printed pair
\((U_{47},U_{111})\) adds only the specialization in which the first two displayed
rows are supplied by \((H_{\mathcal W})\); it is not an additional premise of the
core hierarchy theorem.
\end{corollary}

\begin{proof}
After deleting \((H_{\mathcal W})\), the proof of
\Cref{thm:unconditional-submission-richardson-closure} still gives exactly the
displayed verified-triple obstruction, because that proof uses only
\Cref{thm:richardson-punctured-source-problem,lem:richardson-translation,thm:punctured-domain-richardson-reduction}
and the three semantic rows for the chosen triple.  The internally proved row
\(\Abs=|-|\) on \(\RR\setminus\{0\}\) remains available by
\Cref{thm:eml-punctured-richardson-core}; the first two rows may be supplied by
any finite verified witnesses and do not depend on the printed names.
\Cref{thm:richardson-witness-row-rigidity} proves the converse boundary: if the
printed pair is used as those first two witnesses, then the identities
\(U_{47}=\pi\) and \(U_{111}(x)=\sin x\) on \(\RR\) are precisely the missing row
data.  Hence removing \((H_{\mathcal W})\) deletes only that named
specialization, and leaves no residual unconditional statement about the exact
printed roots \(U_{47}\) and \(U_{111}\).  The downstream obstruction spine and
main hierarchy cite the verified-triple Richardson interface through
\Cref{thm:unconditional-submission-richardson-closure}, so their Richardson input
survives unchanged as the displayed theorem scheme.
\end{proof}

\begin{remark}[Dependency cut]
The source-locking and \(\mathcal B_{36}\)-noninstantiation clauses are separate
from the Richardson obstruction: \Cref{thm:unconditional-submission-richardson-closure} uses
a punctured triple interface and does not import the endpoint-root string.
\end{remark}

\begin{corollary}[Triple-relative expression/protocol separation]
\label{cor:strict-L0-L1}
For every finite punctured EML Richardson triple
\((\mathsf P,\mathsf S,\mathsf A)\) satisfying the three identities of
\Cref{thm:separation}, the imported \(L0\) expression-universality result for
\((\eml,1)\) does not provide a computable semantic normalizer that reflects
and preserves source-domain equality for the translated Richardson fragment.
\end{corollary}

\begin{proof}
A computable semantic normalizer for the translated punctured fragment with the
stated equality-reflection and preservation property would contradict the
normalizer obstruction in \Cref{thm:separation}.  The only
absolute-value identity needed is the punctured row \(|x|\) off zero; no endpoint
square-root identity is used.
\end{proof}

\begin{proposition}[Protocol compilation boundary]
\label{prop:L1-compiles-L0}
Let \(\mathcal E\) be an effective expression grammar with semantic equivalence
relation \(\equiv\), interpreted for a source-indexed partial-real fragment as
equality on the specified common source domains.  Suppose an \(L1\) protocol has
computable normal forms with decidable equality.  If there is a computable
semantic compiler \(C\colon\mathcal E\to N\) into those normal forms such that
\[
  E\equiv E' \quad\Longleftrightarrow\quad C(E)=C(E')
\]
for all source expressions in the fragment under consideration, then semantic
equivalence in \(\mathcal E\) is decidable.  Consequently no Richardson
identity-to-zero source fragment, transported with its domains and retaining this
equality-reflection clause, admits such a compiler into the Zeckendorf \(L1\)
protocol.
\end{proposition}

\begin{proof}
Given two source expressions \(E,E'\), compute \(C(E)
\) and \(C(E')\).  Since equality of target normal forms is decidable in an
\(L1\) protocol, decide whether \(C(E)=C(E')\).  The displayed reflection and
preservation condition makes this decision equivalent to deciding
\(E\equiv E'\), including equality on the specified common domains in the
partial-real case.  If the source fragment is Richardson-undecidable in this
identity-to-zero sense, this would contradict \Cref{lem:richardson-fragment}
after the effective transported translation of \Cref{lem:richardson-translation}.  The Zeckendorf target has
decidable normal-form equality by \Cref{prop:fold-confluence} and is an \(L1\)
protocol by \Cref{thm:zeckendorf-protocol-univ}, so it is covered by the general
argument.
\end{proof}

\begin{corollary}[Compilation scope]
\label{cor:compile-scope}
The Zeckendorf \(L1\) protocol is a computable protocol for natural-number
arithmetic after value transport, but it is not a computable semantic compiler
or equality-reflecting normalizer for an \(L0\) expression grammar whose
source-indexed partial-real equivalence is Richardson-undecidable.
\end{corollary}

\begin{proof}
The positive protocol scope is \Cref{thm:zeckendorf-protocol-univ}.  The excluded
compiler scope is exactly \Cref{prop:L1-compiles-L0}: a compiler reflecting
source-domain semantic equivalence into decidable Zeckendorf normal-form equality
would decide the Richardson-undecidable source identity-to-zero relation.
\end{proof}

\begin{remark}[Transported Richardson boundary]
The Richardson obstruction concerns transported real semantics on the verified
Richardson fragment only.  It is not a decidability statement for arbitrary raw
EML terms on complex principal-branch domains.
Every equality used in \Cref{lem:richardson-translation} is an equality of
partial real functions on the domains of the source Richardson expressions.
The proof never constructs or compares arbitrary EML terms outside that
translated fragment.
\end{remark}

\begin{proposition}[External universality versus verified obstruction]
\label{prop:external-versus-locked-package}
\label{thm:no-hidden-L0-semantic-normalizer}
The Odrzywo{\l}ek \(L0\) import and the Richardson obstruction have different
logical roles: the former supplies expression representatives for a calculator
basis, while the latter requires the explicit three-witness interface of
\Cref{thm:no-core-only-richardson-translation}.
\end{proposition}

\begin{proof}
The \(L0\) import is \Cref{thm:eml-universality}; the obstruction is
\Cref{thm:separation}.  The proof of \Cref{thm:separation} cites only the core
EML witnesses, the three-witness interface, and Richardson's theorem.
\end{proof}

